% problem_id: S001-more-algebra-lemma-special-case
% source_id: S001 (Stacks Project)
% source_file: more-algebra.tex
% statement_locator: more-algebra.tex lines 35369--35383
% proof_locator: more-algebra.tex lines 35385--35511
% env: lemma
% id_note: label 跨文件重名 ⇒ id 用文件限定形式
% pairing: tight (verified)
% status: complete (statement + author proof verbatim + reason + locators)
%
\begin{lemma}
\label{lemma-special-case}
Let $A \subset B \subset C$ be extensions of discrete valuation rings
with fraction fields $K \subset L \subset M$. Assume
\begin{enumerate}
\item the residue field $k$ of $A$ is algebraically closed of
characteristic $p > 0$,
\item $A$ and $B$ are complete,
\item $A \to B$ is weakly unramified,
\item $M$ is a finite extension of $L$,
\item $k = \bigcap\nolimits_{n \geq 1} \kappa_B^{p^n}$
\end{enumerate}
Then there exists a finite extension $K_1/K$ which
is a weak solution for $A \to C$.
\end{lemma}

\begin{proof}
Let $M'$ be any finite extension of $L$ and consider the integral closure
$C'$ of $B$ in $M'$. Then $C'$ is finite over $B$ as $B$ is Nagata by
Algebra, Lemma \ref{algebra-lemma-Noetherian-complete-local-Nagata}.
Moreover, $C'$ is a discrete valuation ring, see discussion in
Remark \ref{remark-construction}. Moreover $C'$ is complete as a
$B$-module, hence complete as a discrete valuation ring, see
Algebra, Section \ref{algebra-section-completion}.
It follows in particular that $C$ is the integral
closure of $B$ in $M$ (by definition of valuation rings as maximal
for the relation of domination).

\medskip\noindent
Let $M \subset M'$ be a finite extension and let $C' \subset M'$
be the integral closure of $B$ as above. By
Lemma \ref{lemma-solutions-go-down}
it suffices to prove the result for $A \to B \to C'$.
Hence we may assume that $M/L$ is normal, see
Fields, Lemma \ref{fields-lemma-normal-closure}.

\medskip\noindent
If $M / L$ is normal, we can find a chain of finite extensions
$$
L = L^0 \subset L^1 \subset L^2 \subset \ldots \subset L^r = M
$$
such that each extension $L^{j + 1}/L^j$ is either:
\begin{enumerate}
\item[(a)] purely inseparable of degree $p$,
\item[(b)] totally ramified with respect to $B^j$ and Galois of degree $p$,
\item[(c)] totally ramified with respect to $B^j$ and Galois cyclic of
order prime to $p$,
\item[(d)] Galois and unramified with respect to $B^j$.
\end{enumerate}
Here $B^j$ is the integral closure of $B$ in $L^j$.
Namely, since $M/L$ is normal we can write it as a compositum of
a Galois extension and a purely inseparable extension
(Fields, Lemma \ref{fields-lemma-normal-case}).
For the purely inseparable extension the existence of the filtration
is clear. In the Galois case, note that $G$ is ``the'' decomposition group
and let $I \subset G$ be the inertia group. Then on the one hand
$I$ is solvable by Lemma \ref{lemma-galois-inertia} and on the other
hand the extension $M^I/L$ is unramified with respect to $B$ by
Lemma \ref{lemma-inertial-invariants-unramified}.
This proves we have a filtration as stated.

\medskip\noindent
We are going to argue by induction on the integer $r$. Suppose that we
can find a finite extension $K_1/K$ which is a weak solution
for $A \to B^1$ where $B^1$ is the integral closure of $B$ in $L^1$.
Let $K'_1$ be the normal closure of $K_1/K$
(Fields, Lemma \ref{fields-lemma-normal-closure}).
Since $A$ is complete and the residue field of $A$ is algebraically closed
we see that $K'_1/K_1$ is separable and totally ramified with
respect to $A_1$ (some details omitted).
Hence $K'_1/K$ is a weak solution for $A \to B^1$ as well by
Lemma \ref{lemma-weakly-unramified-goes-up-along-totally-ramified}.
In other words, we may and do assume that $K_1$ is a normal extension of $K$.
Having done so we consider the sequence
$$
L^0_1 = (L^0 \otimes_K K_1)_{red} \subset
L^1_1 = (L^1 \otimes_K K_1)_{red} \subset \ldots \subset
L^r_1 = (L^r \otimes_K K_1)_{red}
$$
and the corresponding integral closures $B^i_1$. Note that $C_1 = B^r_1$
is a product of discrete valuation rings which are transitively permuted
by $G = \text{Aut}(K_1/K)$ by Lemma \ref{lemma-galois-relative}.
In particular all the extensions of discrete valuation rings
$A_1 \to (C_1)_\mathfrak m$ are isomorphic and a weak solution for one
will be a weak solution for all of them. We can apply the induction
hypothesis to the sequence
$$
A_1 \to (B^1_1)_{B^1_1 \cap \mathfrak m} \to
(B^2_1)_{B^2_1 \cap \mathfrak m} \to
\ldots \to
(B^r_1)_{B^r_1 \cap \mathfrak m} =
(C_1)_\mathfrak m
$$
to get a weak solution $K_2/K_1$ for $A_1 \to (C_1)_\mathfrak m$.
The extension $K_2/K$ will then be a weak solution for $A \to C$
by what we said before. Note that the induction hypothesis applies:
the ring map $A_1 \to (B^1_1)_{B^1_1 \cap \mathfrak m}$
is weakly unramified by our choice of $K_1$
and the sequence of fraction field extensions
each still have one of the properties (a), (b), (c), or (d)
listed above. Moreover, observe that for any finite extension 
$\kappa_B \subset \kappa$ we still have $k = \bigcap \kappa^{p^n}$.

\medskip\noindent
Thus everything boils down to finding a weak solution for $A \subset C$
when the field extension $M/L$ satisfies one of the properties
(a), (b), (c), or (d).

\medskip\noindent
Case (d). This case is trivial as here $B \to C$ is unramified already.

\medskip\noindent
Case (c). Say $M/L$ is cyclic of order $n$ prime to $p$. Because
$M/L$ is totally ramified with respect to $B$ we see that the ramification
index of $B \subset C$ is $n$ and hence the ramification index of $A \subset C$
is $n$ as well. Choose a uniformizer $\pi \in A$ and set
$K_1 = K[\pi^{1/n}]$. Then $K_1/K$ is a solution for $A \subset C$
by Abhyankar's lemma (Lemma \ref{lemma-abhyankar}).

\medskip\noindent
Case (b). We divide this case into the mixed characteristic case and the
equicharacteristic case. In the equicharacteristic case this is
Lemma \ref{lemma-characteristic-p-case}. In the mixed characteristic
case, we first replace $K$ by a finite extension to get to the
situation where $M/L$ is a degree $p$ extension of finite level using
Lemma \ref{lemma-make-finite-level}.
Then the level is a rational number $l \in [0, p)$, see discussion
preceding Lemma \ref{lemma-lowering-the-level}. If the level is $0$,
then $B \to C$ is weakly unramified and we're done. If not, then we
can replacing the field $K$ by a finite extension to obtain a new
situation with level $l' \leq \max(0, l - 1, 2l - p)$ by
Lemma \ref{lemma-lowering-the-level}.
If $l = p - \epsilon$ for $\epsilon < 1$ then we see that
$l' \leq p - 2\epsilon$. Hence after a finite number of replacements
we obtain a case with level $\leq p - 1$. Then after at most $p - 1$
more such replacements we reach the situation where the level is zero.

\medskip\noindent
Case (a) is Lemma \ref{lemma-purely-inseparable-case}. This is the only case
where we possibly need a purely inseparable extension of $K$, namely, in
case (2) of the statement of the lemma we win by adjoining a $p$th power
of the element $\pi$. This finishes the proof of the lemma.
\end{proof}
